\documentclass[12pt,a4paper]{article}
\usepackage[utf8]{inputenc}
\usepackage[vietnamese]{babel}
\usepackage{amsmath,amssymb,amsthm}
\usepackage{geometry}
\geometry{left=2cm,right=2cm,top=2cm,bottom=2cm}
\usepackage[most]{tcolorbox}
\usepackage{hyperref}

\title{\textbf{HƯỚNG DẪN GIẢI CHI TIẾT CÁC DẠNG BÀI TẬP\\PHÉP BIẾN HÌNH PHÂN TUYẾN TÍNH \& ĐỐI XỨNG QUA ĐƯỜNG TRÒN}}
\author{Chuyên đề Hàm Biến Phức}
\date{\today}

\tcbset{colback=blue!5!white,colframe=blue!75!black,fonttitle=\bfseries,arc=2mm}

\begin{document}

\maketitle

\tableofcontents
\newpage

% ==============================================================================
\section{Dạng 1: Tìm ảnh và tạo ảnh qua phép biến hình phân tuyến tính}
% ==============================================================================

\begin{tcolorbox}[title=Phương pháp giải tổng quát]
Phép biến hình phân tuyến tính có dạng tổng quát:
\[ w = f(z) = \frac{az + b}{cz + d} \quad (ad - bc \neq 0) \]
\begin{enumerate}
    \item \textbf{Tìm ảnh của miền $D_z$ qua phép biến hình $w = f(z)$}:
    \begin{itemize}
        \item Biểu diễn $z$ theo $w$: Rút $z = f^{-1}(w) = \frac{-dw + b}{cw - a}$.
        \item Thay biểu thức $z(w)$ vào bất phương trình/phương trình xác định miền $D_z$.
        \item Sử dụng tọa độ đại số $w = u + iv$ hoặc modun $|w|$ để đưa về phương trình miền $D_w$ trong mặt phẳng $w$.
    \end{itemize}
    
    \item \textbf{Tìm tạo ảnh của miền $D_w$ qua phép biến hình $w = f(z)$}:
    \begin{itemize}
        \item Thay biểu thức $w = \frac{az + b}{cz + d}$ trực tiếp vào bất phương trình/phương trình xác định miền $D_w$.
        \item Đặt $z = x + iy$, sử dụng biến đổi số phức để suy ra điều kiện của $(x, y)$ trong mặt phẳng $z$.
    \end{itemize}
\end{enumerate}
\end{tcolorbox}

\subsection*{Bài tập mẫu: Ví dụ 2.11}
Cho $w = \frac{2z + 1}{z - 1}$. Trong $\overline{\mathbb{C}}$, tìm:
\begin{enumerate}
    \item[a)] Ảnh của $x^2 + y^2 \le 1$.
    \item[b)] Ảnh của $\text{Re}z > 2$.
    \item[c)] Tạo ảnh của $u^2 + v^2 < 1$.
    \item[d)] Tạo ảnh của $\text{Re}w = 2$.
\end{enumerate}

\subsubsection*{Lời giải chi tiết}

Rút $z$ theo $w$:
\[ w(z - 1) = 2z + 1 \iff wz - w = 2z + 1 \iff z(w - 2) = w + 1 \iff z = \frac{w + 1}{w - 2} \]

\noindent \textbf{a) Tìm ảnh của $x^2 + y^2 \le 1$:}

Miền $x^2 + y^2 \le 1$ tương đương với $|z| \le 1$. Thay $z = \frac{w + 1}{w - 2}$ vào bất phương trình:
\[ \left| \frac{w + 1}{w - 2} \right| \le 1 \iff |w + 1| \le |w - 2| \]
Đặt $w = u + iv$ ($u, v \in \mathbb{R}$):
\[ |(u + 1) + iv|^2 \le |(u - 2) + iv|^2 \iff (u + 1)^2 + v^2 \le (u - 2)^2 + v^2 \]
\[ \iff u^2 + 2u + 1 + v^2 \le u^2 - 4u + 4 + v^2 \iff 6u \le 3 \iff u \le \frac{1}{2} \]
\textbf{Kết luận:} Ảnh cần tìm là nửa mặt phẳng $\text{Re}(w) \le \frac{1}{2}$.

\bigskip
\noindent \textbf{b) Tìm ảnh của $\text{Re}z > 2$:}

Ta có $\text{Re}z = \frac{z + \bar{z}}{2} > 2 \iff z + \bar{z} > 4$. Thay $z = \frac{w + 1}{w - 2}$:
\[ \frac{w + 1}{w - 2} + \frac{\bar{w} + 1}{\bar{w} - 2} > 4 \iff \frac{(w + 1)(\bar{w} - 2) + (\bar{w} + 1)(w - 2)}{|w - 2|^2} > 4 \]
Khai triển tử số:
\[ (w\bar{w} - 2w + \bar{w} - 2) + (w\bar{w} - 2\bar{w} + w - 2) = 2|w|^2 - (w + \bar{w}) - 4 = 2(u^2 + v^2) - 2u - 4 \]
Bất phương trình trở thành:
\[ \frac{2(u^2 + v^2) - 2u - 4}{(u - 2)^2 + v^2} > 4 \iff 2u^2 + 2v^2 - 2u - 4 > 4(u^2 - 4u + 4 + v^2) \]
\[ \iff 2u^2 - 14u + 2v^2 + 20 < 0 \iff u^2 - 7u + v^2 + 10 < 0 \]
\[ \iff \left(u - \frac{7}{2}\right)^2 + v^2 < \frac{49}{4} - 10 = \frac{9}{4} = \left(\frac{3}{2}\right)^2 \]
\textbf{Kết luận:} Ảnh cần tìm là hình tròn mở tâm $w_0 = \frac{7}{2}$, bán kính $R = \frac{3}{2}$, tức là $\left| w - \frac{7}{2} \right| < \frac{3}{2}$.

\bigskip
\noindent \textbf{c) Tìm tạo ảnh của $u^2 + v^2 < 1$:}

Miền $u^2 + v^2 < 1$ tương đương $|w| < 1$. Thay $w = \frac{2z + 1}{z - 1}$ vào:
\[ \left| \frac{2z + 1}{z - 1} \right| < 1 \iff |2z + 1| < |z - 1| \]
Đặt $z = x + iy$ ($x, y \in \mathbb{R}$):
\[ |(2x + 1) + 2iy|^2 < |(x - 1) + iy|^2 \iff (2x + 1)^2 + 4y^2 < (x - 1)^2 + y^2 \]
\[ \iff 4x^2 + 4x + 1 + 4y^2 < x^2 - 2x + 1 + y^2 \iff 3x^2 + 6x + 3y^2 < 0 \]
\[ \iff x^2 + 2x + y^2 < 0 \iff (x + 1)^2 + y^2 < 1 \]
\textbf{Kết luận:} Tạo ảnh cần tìm là hình tròn mở tâm $z_0 = -1$, bán kính $R = 1$, tức là $|z + 1| < 1$.

\bigskip
\noindent \textbf{d) Tìm tạo ảnh của $\text{Re}w = 2$:}

Đường $\text{Re}w = 2 \iff \frac{w + \bar{w}}{2} = 2 \iff w + \bar{w} = 4$. Thay $w = \frac{2z + 1}{z - 1}$:
\[ \frac{2z + 1}{z - 1} + \frac{2\bar{z} + 1}{\bar{z} - 1} = 4 \iff \frac{(2z + 1)(\bar{z} - 1) + (2\bar{z} + 1)(z - 1)}{|z - 1|^2} = 4 \]
Tử số:
\[ (2z\bar{z} - 2z + \bar{z} - 1) + (2z\bar{z} - 2\bar{z} + z - 1) = 4|z|^2 - (z + \bar{z}) - 2 = 4(x^2 + y^2) - 2x - 2 \]
Do đó:
\[ \frac{4(x^2 + y^2) - 2x - 2}{(x - 1)^2 + y^2} = 4 \iff 4x^2 + 4y^2 - 2x - 2 = 4(x^2 - 2x + 1 + y^2) \]
\[ \iff -2x - 2 = -8x + 4 \iff 6x = 6 \iff x = 1 \]
\textbf{Kết luận:} Tạo ảnh cần tìm là đường thẳng $\text{Re}z = 1$ (đường $x = 1$).

\newpage
% ==============================================================================
\section{Dạng 2: Phép đối xứng qua đường tròn}
% ==============================================================================

\begin{tcolorbox}[title=Định nghĩa và Công thức điểm đối xứng qua đường tròn]
Cho đường tròn $S(z_0, r) = \{z \in \mathbb{C} : |z - z_0| = r\}$.
Hai điểm $z$ và $z^*$ được gọi là đối xứng với nhau qua đường tròn $S(z_0, r)$ nếu:
\begin{enumerate}
    \item $z, z_0, z^*$ thẳng hàng và nằm cùng phía đối với tâm $z_0$ (tức là $\arg(z - z_0) = \arg(z^* - z_0)$).
    \item Tích khoảng cách: $|z - z_0| \cdot |z^* - z_0| = r^2$.
\end{enumerate}
\textbf{Công thức tính nhanh điểm đối xứng $z^*$ của $z$:}
\[ z^* - z_0 = \frac{r^2}{\overline{z - z_0}} \iff z^* = z_0 + \frac{r^2}{\overline{z - z_0}} \]
\end{tcolorbox}

\subsection*{Bài tập mẫu: Ví dụ 2.14}
Tìm điểm $z^*$ đối xứng với $z = 1 + i$ qua đường tròn $S_{0,2}$ (đường tròn tâm $z_0 = 0$, bán kính $r = 2$).

\subsubsection*{Lời giải chi tiết}

\noindent \textbf{Cách 1 (Theo định nghĩa):}
Đặt $z_0 = 0, r = 2$.
Vì $\arg(z - z_0) = \arg(z^* - z_0)$ nên tồn tại số thực $k > 0$ sao cho $z^* - z_0 = k(z - z_0) \implies z^* = k z$.
Mặt khác:
\[ |z - z_0| \cdot |z^* - z_0| = r^2 \iff |z| \cdot |k z| = 4 \iff k |z|^2 = 4 \]
Với $z = 1 + i \implies |z|^2 = 1^2 + 1^2 = 2$. Suy ra $2k = 4 \iff k = 2$.
Vậy $z^* = 2z = 2(1 + i) = 2 + 2i$.

\bigskip
\noindent \textbf{Cách 2 (Dùng công thức cực nhanh):}
Áp dụng trực tiếp công thức:
\[ z^* = z_0 + \frac{r^2}{\overline{z - z_0}} = 0 + \frac{2^2}{\overline{1 + i}} = \frac{4}{1 - i} = \frac{4(1 + i)}{(1 - i)(1 + i)} = \frac{4(1 + i)}{2} = 2 + 2i \]

\subsection*{Bài tập mẫu: Ví dụ 2.15}
Tìm ảnh bởi phép đối xứng qua đường tròn $|z - 2| = 2$ của:
\begin{enumerate}
    \item[a)] $|z - 2| = \frac{1}{2}$
    \item[b)] $|z - 1| = 2$
    \item[c)] $\text{Im}z = 2$
\end{enumerate}

\subsubsection*{Lời giải chi tiết}

Đường tròn đối xứng có tâm $z_0 = 2$, bán kính $r = 2$.
Công thức phép đối xứng biến điểm $z$ thành $z^*$:
\[ z^* - 2 = \frac{4}{\bar{z} - 2} \implies \bar{z} - 2 = \frac{4}{z^* - 2} \implies z = 2 + \frac{4}{\bar{z}^* - 2} \]

\noindent \textbf{a) Ảnh của $|z - 2| = \frac{1}{2}$:}

Thay $z - 2 = \frac{4}{\bar{z}^* - 2}$ vào phương trình:
\[ \left| \frac{4}{\bar{z}^* - 2} \right| = \frac{1}{2} \iff |\bar{z}^* - 2| = 8 \iff |z^* - 2| = 8 \]
\textbf{Kết luận:} Ảnh là đường tròn tâm $z_0 = 2$, bán kính $R = 8$.

\bigskip
\noindent \textbf{b) Ảnh của $|z - 1| = 2$:}

Thay $z = 2 + \frac{4}{\bar{z}^* - 2}$ vào:
\[ \left| 1 + \frac{4}{\bar{z}^* - 2} \right| = 2 \iff \left| \frac{\bar{z}^* + 2}{\bar{z}^* - 2} \right| = 2 \iff |\bar{z}^* + 2| = 2 |\bar{z}^* - 2| \]
Bình phương 2 vế với $z^* = u + iv$:
\[ (u + 2)^2 + v^2 = 4[(u - 2)^2 + v^2] \iff u^2 + 4u + 4 + v^2 = 4(u^2 - 4u + 4 + v^2) \]
\[ \iff 3u^2 - 20u + 3v^2 + 12 = 0 \iff u^2 - \frac{20}{3}u + v^2 + 4 = 0 \]
\[ \iff \left(u - \frac{10}{3}\right)^2 + v^2 = \frac{100}{9} - 4 = \frac{64}{9} = \left(\frac{8}{3}\right)^2 \]
\textbf{Kết luận:} Ảnh là đường tròn tâm $w_0 = \frac{10}{3}$, bán kính $R = \frac{8}{3}$.

\bigskip
\noindent \textbf{c) Ảnh của $\text{Im}z = 2$:}

Ta có $\text{Im}z = \frac{z - \bar{z}}{2i} = 2 \iff z - \bar{z} = 4i$. Thay $z = 2 + \frac{4}{\bar{z}^* - 2}$ và $\bar{z} = 2 + \frac{4}{z^* - 2}$:
\[ \frac{4}{\bar{z}^* - 2} - \frac{4}{z^* - 2} = 4i \iff \frac{1}{\bar{z}^* - 2} - \frac{1}{z^* - 2} = i \iff \frac{(z^* - 2) - (\bar{z}^* - 2)}{|z^* - 2|^2} = i \]
\[ \iff \frac{z^* - \bar{z}^*}{|z^* - 2|^2} = i \iff \frac{2i \text{Im}(z^*)}{|z^* - 2|^2} = i \iff \frac{2v}{(u - 2)^2 + v^2} = 1 \]
\[ \iff 2v = (u - 2)^2 + v^2 \iff (u - 2)^2 + v^2 - 2v = 0 \iff (u - 2)^2 + (v - 1)^2 = 1 \]
\textbf{Kết luận:} Ảnh là đường tròn $(u - 2)^2 + (v - 1)^2 = 1$ (tâm $(2, 1)$, bán kính $R = 1$), trừ điểm $(2, 0)$ vì điểm vô cực biến thành $z^* = 2$.

\newpage
% ==============================================================================
\section{Dạng 3: Tìm phép biến hình phân tuyến tính biến miền này thành miền khác}
% ==============================================================================

\begin{tcolorbox}[title=Phương pháp tìm phép biến hình phân tuyến tính]
Để tìm phép biến hình phân tuyến tính $w = f(z)$ biến miền $D_1$ thành miền $D_2$:
\begin{enumerate}
    \item \textbf{Xác định nghiệm của $w$ và cực của $w$}:
    Lấy điểm $z_0 \in D_1$ sao cho $w(z_0) = 0 \in D_2$.
    \item \textbf{Sử dụng nguyên lý bảo toàn tính đối xứng}:
    \begin{itemize}
        \item Tìm điểm $z_0^*$ đối xứng với $z_0$ qua biên $\partial D_1$.
        \item Tìm điểm $w_0^*$ đối xứng với $w(z_0) = 0$ qua biên $\partial D_2$.
        \item Vì phép biến hình phân tuyến tính bảo toàn tính đối xứng, ta có $w(z_0^*) = w_0^* = \infty$.
    \end{itemize}
    \item \textbf{Thiết lập dạng hàm $w(z)$}:
    \[ w(z) = k \cdot \frac{z - z_0}{z - z_0^*} \quad (k \in \mathbb{C}, k \neq 0) \]
    \item \textbf{Xác định hằng số $k$}:
    Lấy điểm $z$ thuộc biên $\partial D_1$ thì $w(z)$ thuộc biên $\partial D_2$. Sử dụng điều kiện biên để tìm Hằng số $|k|$.
\end{enumerate}
\end{tcolorbox}

\subsection*{Bài tập mẫu: Ví dụ 2.16}
Tìm tất cả các phép biến hình phân tuyến tính biến nửa mặt phẳng $\text{Im}z > 0$ thành hình tròn đơn vị $|w| < 1$.

\subsubsection*{Lời giải chi tiết}

Giả sử $w = f(z)$ là phép biến hình cần tìm.
\begin{enumerate}
    \item Đặt $w^{-1}(0) = z_0$, nghĩa là $w(z_0) = 0$. Vì $0 \in \{|w| < 1\}$ nên điểm $z_0$ phải thuộc nửa mặt phẳng trên, tức là $\text{Im}z_0 > 0$.
    
    \item Xác định các điểm đối xứng:
    \begin{itemize}
        \item Biên của miền gốc $\text{Im}z > 0$ là trục thực $\text{Im}z = 0$. Điểm đối xứng với $z_0$ qua trục thực là $\bar{z}_0$.
        \item Biên của miền đích $|w| < 1$ là đường tròn đơn vị $|w| = 1$. Điểm đối xứng với $0$ qua đường tròn $|w| = 1$ là điểm vô cực $\infty$.
    \end{itemize}
    Do phép biến hình phân tuyến tính bảo toàn tính đối xứng nên điểm $\bar{z}_0$ phải biến thành điểm $\infty$, tức là $w(\bar{z}_0) = \infty$.
    
    \item Do đó, hàm $w(z)$ có dạng:
    \[ w(z) = k \cdot \frac{z - z_0}{z - \bar{z}_0} \quad (k \in \mathbb{C}, k \neq 0) \]
    
    \item Xác định $k$:
    Lấy $x \in \mathbb{R}$ bất kỳ thuộc biên $\text{Im}z = 0$. Khi đó $w(x)$ phải thuộc biên $|w| = 1$, tức là $|w(x)| = 1$:
    \[ |w(x)| = |k| \cdot \frac{|x - z_0|}{|x - \bar{z}_0|} = 1 \]
    Vì $x \in \mathbb{R}$ nên khoảng cách từ $x$ đến $z_0$ và $\bar{z}_0$ bằng nhau: $|x - z_0| = |x - \bar{z}_0|$.
    Do đó $|k| = 1 \implies k = e^{i\theta}$ với $\theta \in \mathbb{R}$.
\end{enumerate}

\textbf{Kết luận:} Tất cả các phép biến hình phân tuyến tính cần tìm có dạng:
\[ w(z) = e^{i\theta} \frac{z - z_0}{z - \bar{z}_0} \quad \text{với } \text{Im}z_0 > 0 \text{ và } \theta \in \mathbb{R} \]

\end{document}
